\clearpage
\section*{Raw permit delays and adjusted local score orderings}
\textit{September 15, 2026. Agent-drafted comparison and manuscript proposal requested by Jacob.}

The score adjustment changes a substantial fraction of the local orderings
used in the paper. Raw arithmetic mean processing days agree with the
score for 33 of 44 permit-event-study pairs (75.0\%), 90 of 161 multifamily
pairs (55.9\%), 222 of 367 sales pairs (60.5\%), and 142 of 244 rental pairs
(58.2\%). For all construction, agreement is 152 of 256 pairs (59.4\%).
The permit comparison uses scores and raw times through 2014; the boundary
outcomes use both measures through 2022. Each comparison uses the same
retained permits as its score: 71,059 observations for 83 aldermen through
2014 and 120,505 observations for 113 aldermen through 2022. No raw-time
tie occurs for the arithmetic means.

The share of observations represented by agreeing pairs differs from the
unweighted pair share. Agreement covers 76.1\% of reassigned permit blocks,
51.4\% of multifamily projects, 57.1\% of sales, and only 39.1\% of rental
observations; the all-construction share is 53.2\%. These weights describe
sample representation, not regression leverage or contributions to an
estimate. The rental gap reflects several large disagreements. Reilly's
raw mean is 37.4 days and Hopkins's is 69.4 days, yet the adjusted score
ranks Reilly as more stringent. That pair represents 21,216 rental
observations. Reilly--Burnett, Hopkins--Burnett, and Cappleman--Tunney also
reverse the raw ordering and together represent another 48,631 rental
observations. These are measured disagreements, not identified data errors.

Part of the difference reflects the outcome scale. Using mean log days,
without regression adjustment, increases agreement to 38 of 44 permit pairs
(86.4\%), 109 of 161 multifamily pairs (67.7\%), 266 of 367 sales pairs
(72.5\%), and 174 of 244 rental pairs (71.3\%). The corresponding
observation shares are 91.7\%, 62.5\%, 70.6\%, and 50.5\%. Thus taking
logs explains some disagreement, but the remaining adjustment still
reverses the ordering for roughly half of rental observations. Median days
and raw ties are also retained in the report. The arithmetic-mean/raw-score
citywide rank correlation is 0.573 through 2014 and 0.286 through 2022;
using mean log days gives 0.622 and 0.425. All-citywide comparisons include
aldermen who did not serve together and are contextual rather than the
paper's identifying comparisons.

This exercise describes what the score construction changes. Both measures
use the same permits, so agreement is not independent validation of
alderman behavior. Raw averages combine permit mix, calendar time, location
and ward characteristics. The adjusted score also changes aggregation and
applies shrinkage. This comparison does not isolate the contribution of
each control or determine which ordering is closer to a causal effect.
The more closely matched mean-log comparison removes the transformation
difference while preserving the other differences.

Jacob also asked about adding score-robustness checks to the paper. The
review proposes a short main-text pointer and an appendix section on local
ordering stability and its consequences for the outcome estimates. A
concrete five-outcome table reports baseline estimates, score-draw ranges,
sign retention, deletion estimates, conditional p-values and retained
sample shares. All three reweighting methods and both 75\% and 95\%
thresholds are preserved, including unfavorable reversal tests. The proposed
text distinguishes the comparatively stable permit and price signs from
the more sensitive multifamily DUPAC estimate and the imprecise rental
estimate. It does not establish persistence over time or substitute for the
earlier weak chronological split. The proposal also suggests removing
``persistent'' from the current Appendix B description of the alderman
effect. No manuscript changes have been applied.

\clearpage
\begin{center}
\includegraphics[width=\textwidth]{../tasks/audits/score_robustness_review/output/raw_processing_comparison.pdf}
\end{center}

\textit{Reproduction and definitions.} Run \texttt{make} in
\path{tasks/audits/score_robustness_review/code/}. The audit reads the
preserved permit source under \path{data_raw/score_robustness/}; the source
hash and producer are documented in the score-reliability audit. The raw
comparison reproduces that audit's eligibility rules, including lagged
workload and complete controls, and verifies each alderman's retained
permit and observed-month counts against the saved baseline. Raw averages
weight each retained permit equally. Same-day permits remain in workload
counts but do not enter either processing-time measure. Dates and source
processing times are not changed or newly trimmed.

One script calculates raw alderman averages, all local and citywide pair
comparisons, agreement counts, correlations and the figure. The second
assembles the existing outcome-sensitivity results and proposed wording
into a review HTML. Outputs contain 196 alderman-period rows, 32,409
pair/period/measure/sample rows, 21 agreement summaries, six correlations,
and 30 combined robustness rows. Keys, raw-time/log consistency, score
coverage, and agreement/reversal/tie totals are checked. A raw tie remains
in the denominator as neither agreement nor reversal. Formal rank
confidence sets require additional inference beyond the perturbation
frequencies reported here; the review links to Mogstad, Romano, Shaikh and
Wilhelm (2024) to distinguish those concepts. Production inputs, score
estimation, outcome estimates and manuscript files are unchanged.
